\section{Kepler problem and Poincar\'e coordinates}\label{bpP:p7}
\begingroup\sloppy

The unperturbed part of both physical models is a sum of $n$ decoupled planar Kepler problems (Lemma~\ref{bpP:p6:split}). This part
gives the coordinates in which that sum is integrable \emph{and} regular at the circular orbits: the Poincar\'e chart
$(\lambda,\Lambda,\xi,\eta)$. The results are:
\begin{itemize}
\item \texttt{kepler} (Prop.~\ref{bpP:p7:kepler}): the planar Kepler problem of a planet of mass $m$ in the field of a centre of strength $\kappa=Gm_0$, through the
mass/strength dictionary of \Iref{bpI:p7:scaling}: the Delaunay map with mass is symplectic (factor one) in the variables $(\ell,g,\Lambda,C)$ with $\Lambda=m\sqrt{\kappa a}$,
$H_{\rm Kep}=-\kappa^2m^3/(2\Lambda^2)$, and the mean values $\langle1/\abs q\rangle$, $\langle\abs p^2\rangle$, $\langle\abs q\rangle$;
\item \texttt{regular-kepler} (Prop.~\ref{bpP:p7:regular-kepler}): Kepler's equation in the eccentric longitude $F-v_0\sin F+v_1\cos F=\lambda$, with the eccentricity \emph{vector} $v$ as parameter, solved
analytically (real, and holomorphically near $v=0$), with no angle of perihelion;
\item \texttt{poincare} (Thm.~\ref{bpP:p7:poincare}): the Poincar\'e chart, real analytic through $\xi=\eta=0$ (circular orbits), a $1$-chart in the sense of
Definition~\ref{bpP:p1:chart}, bijective modulo $2\pi$ in $\lambda$ onto the prograde bound region, with the inverse formulas; Prop.~\ref{bpP:p7:holomorphic} is its holomorphic
extension near the circular torus;
\item \texttt{product-chart} (Thm.~\ref{bpP:p7:product-chart}): $n$ independent charts; the \emph{same} chart serves $H_{\rm mov}$ and $H_{\rm fix}$, with $H^\Psi=h(\Lambda)+\mu f_{\rm mov}\circ\Psi$,
$h(\Lambda)+\varepsilon f_{\rm fix}\circ\Psi$; the tube lemma, and the lattice hypotheses (Q1), (Q2) of Theorem~\ref{bpP:p1:torus-transport} (Lemmas~\ref{bpP:p7:tube}, \ref{bpP:p7:lattice});
\item \texttt{eccentric-longitude} (Prop.~\ref{bpP:p7:d25}): the D25 secular definitions (\texttt{ellipsePosition}, \texttt{ellipseMomentum}, \texttt{keplerMeanLongitude}, \texttt{longitudeJacobian}, \texttt{secularEccVector}, \dots)
are exactly the chart in the eccentric longitude $F$; and the change of variable $\lambda\leftrightarrow F$ in angular averages (Cor.~\ref{bpP:p7:average}).
\end{itemize}
\emph{Ownership.} The Kepler fundamentals are owned by track~I Part~7 and are cited, not repeated: Kepler's equation on a complex tube around the real eccentricity interval, including $e=0$ (\Iref{bpI:p7:kepler-eq}),
the Delaunay orbit maps (\Iref{bpI:p7:orbit-maps}, \Iref{bpI:p7:orbit-maps-e}), their identities (\Iref{bpI:p7:orbit-identities}), the circular endpoint (\Iref{bpI:p7:circular}), the mean values (\Iref{bpI:p7:means}), the dictionary
(\Iref{bpI:p7:scaling}), and the Delaunay chart (\Iref{bpI:p7:symplectic}, \Iref{bpI:p7:chart}). Part~7 owns the Poincar\'e chart and the product chart.

\subsection{Conventions, external facts, notation}\label{bpP:p7:conv}

Throughout $m>0$ is the planet mass and $\kappa=Gm_0>0$ the strength of the centre; $\pl{\cdot}$ is the planar length (Definition~\ref{bpP:p0:plane}); angles $\lambda,\ell,g,F$ are in \emph{radians} (period $2\pi$), as in the D25 secular
definitions; the passage to the frozen period-one torus parameter is Lemma~\ref{bpP:p1:period-convert} and is made in Lemma~\ref{bpP:p7:lattice}. The frozen sign convention is that of Definition~\ref{bpP:p0:hamsymp}: a phase point is
$z=(z_1,z_2)$ with positions first, and $\Omega=\sum dz_{2,i}\wedge dz_{1,i}$. In the Poincar\'e coordinates of one planet $(\lambda,\xi;\Lambda,\eta)$ the positions are $(\lambda,\xi)$ and the momenta $(\Lambda,\eta)$
(this is the order of \Lean{secularGamma}: $\Gamma=(\xi^2+\eta^2)/2=R(\eta,\xi)$, Definition~\ref{bpP:p0:oscaction}), so
\[\Omega_{\rm chart}=d\Lambda\wedge d\lambda+d\eta\wedge d\xi .\]
Write $J_0v=(-v_1,v_0)$ (\Lean{quarterTurn}), $\mathsf R(g)$ for the rotation by $g$, $e_g=\mathsf R(g)(1,0)=(\cos g,\sin g)$, and $\operatorname{cr}(q,p)=q_0p_1-q_1p_0$ (\Lean{planeCross}).
The Delaunay map of track~I (strength one, unit mass) is written $\mathsf Q^{\rm I}(\ell,g,L,G)$ (\Iref{bpI:p7:orbit-maps}); it satisfies $(\mathsf Q^{\rm I})^*\bigl(\sum dq_i\wedge dp_i\bigr)=d\ell\wedge dL+dg\wedge dG$ (\Iref{bpI:p7:symplectic}).

\emph{External facts.}
\begin{enumerate}[label=(E\arabic*),ref=(E\arabic*)]
\item\label{bpP:p7:e-ift} \emph{Real-analytic inverse function theorem} (as \ref{bpP:p1:ext-ift}, with analytic regularity): a real-analytic map with invertible differential at $z_0$ is a real-analytic diffeomorphism from a neighbourhood of $z_0$ onto an open set; in particular an
implicit equation $\Phi(F,x)=0$ with $\partial_F\Phi\ne0$ has a locally unique real-analytic solution $F(x)$. A bijective real-analytic map between open sets with everywhere invertible differential is a real-analytic diffeomorphism.
\item\label{bpP:p7:e-hol} \emph{Holomorphic calculus.} Sums, products, compositions and quotients (nonvanishing denominator) of holomorphic functions of several variables are holomorphic; $\sqrt{\cdot}$ is holomorphic on $\{\operatorname{Re}>0\}$; a locally uniform limit of holomorphic functions is holomorphic (\Iref{bpI:p7:x-hol}, \Iref{bpI:p7:x-weier}).
\item\label{bpP:p7:e-banach} Banach's fixed point theorem, the mean value inequality on a convex set (\Kref{bp:p1:ext-banach}, \Kref{bp:p1:ext-mvt}).
\item\label{bpP:p7:e-forms} \emph{Exterior calculus on open sets of $\R^N$} (\Iref{bpI:p7:x-forms}): $d$ commutes with pull-back; the identity $\Omega(Au,Av)=c\,\Omega(u,v)$ for the differential $A=D\Psi$ holds iff the pulled-back $2$-forms satisfy $\Psi^*\Omega=c\,\Omega$ (Lemma~\ref{bpP:p1:omega-algebra}(c)).
\item\label{bpP:p7:e-dense} A continuous identity between continuous functions which holds on a dense subset of an open set holds on the whole set.
\item\label{bpP:p7:e-cov} \emph{Change of variables in $\R$} (\Iref{bpI:p7:x-change}).
\item\label{bpP:p7:e-comp} Compact sets and open sets: a compact $K$ contained in an open $V\subset\R^N$ has a positive max-norm distance to $\R^N\setminus V$.
\end{enumerate}

\subsection{The planar Kepler problem with mass $m$ and strength $\kappa$}\label{bpP:p7:sec-kepler}

\begin{definition}[Kepler data]\label{bpP:p7:kepler-data}
For $(q,p)\in\R^2\times\R^2$, $q\ne0$, put
\[E_{m,\kappa}(q,p)=\frac{\pl p^2}{2m}-\frac{\kappa m}{\pl q},\qquad C(q,p)=\operatorname{cr}(q,p),\qquad a=-\frac{\kappa m}{2E_{m,\kappa}},\qquad\Lambda=m\sqrt{\kappa a}\ \ (E_{m,\kappa}<0).\]
The \emph{prograde bound region} is $\mathcal K_{m,\kappa}=\{(q,p):q\ne0,\ E_{m,\kappa}<0,\ C>0\}$ (open: strict inequalities between continuous functions on $\{q\ne0\}$). On it the eccentricity is $e=\sqrt{1-C^2/\Lambda^2}$
(Lemma~\ref{bpP:p7:range}: $0<C\le\Lambda$). The \emph{Delaunay map with mass} is defined on $\mathcal D^m=\{(\ell,g,\Lambda,C)\in\R^4:\Lambda>C>0\}$ by
\[\mathsf Q_{m,\kappa}(\ell,g,\Lambda,C)=\mathsf S_{m,\kappa}\circ\mathsf Q^{\rm I}\Bigl(\ell,g,\frac\Lambda m,\frac Cm\Bigr),\qquad\mathsf S_{m,\kappa}(Q,P)=\Bigl(\frac Q\kappa,\ m\kappa P\Bigr).\]
For $\kappa=m=1$ this is $\mathsf Q^{\rm I}$. For $E_{m,\kappa}<0$ and $m=1$ the quantity $\Lambda=\sqrt{\kappa a}$ and the orbit is the one of $E_\kappa$ in \Iref{bpI:p7:scaling}.
\frozen{KAM.Planetary.keplerEnergy, keplerAxis, keplerAction, planeCross (the case $\kappa=1$, $n$ planets componentwise)}
\source{Fej04r (31)--(32), p.~47--48 ($\Lambda_j=\mu_j\sqrt{M_ja_j}$, $F_{\rm Kep}=-\mu_j^3M_j^2/(2\Lambda_j^2)$; here the mass $m_j$ and strength $\kappa$ replace the reduced mass and total mass, see the deviations); \Iref{bpI:p7:scaling} for $\mathsf S_{m,\kappa}$ (this is its mass dictionary).}
\end{definition}

\begin{lemma}[Range of the elements]\label{bpP:p7:range}
For $(q,p)\in\mathcal K_{m,\kappa}$: $0<C\le\Lambda$, with equality iff the orbit is circular ($\pl q$ constant). Equivalently $\Gamma:=\Lambda-C\in[0,\Lambda)$.
\end{lemma}
\begin{proof}
\proofstep{Scaling to strength one}\label{stepP:p7:rg-scale} With $(Q,P)=\mathsf S_{m,\kappa}^{-1}(q,p)=(\kappa q,p/(m\kappa))$: $E_1(Q,P)=\frac1{m\kappa^2}E_{m,\kappa}(q,p)$ (\Iref{bpI:p7:scaling}: $E_\kappa\circ\mathsf S_\kappa=\kappa^2E_1$ and $E_{m,\kappa}=mE_\kappa(q,p/m)$), $\operatorname{cr}(Q,P)=C/m$, and $L:=\Lambda/m=\sqrt{-1/(2E_1(Q,P))}$.
\proofstep{Lagrange identity}\label{stepP:p7:rg-lagrange} With $r=\pl Q$, $\pl Q^2\pl P^2=(Q\cdot P)^2+\operatorname{cr}(Q,P)^2$ and $\frac12\pl P^2=E_1+1/r=-\frac1{2L^2}+\frac1r$. Hence
$\operatorname{cr}(Q,P)^2=2r-r^2/L^2-(Q\cdot P)^2\le2r-r^2/L^2=L^2-(r-L^2)^2/L^2\le L^2$, i.e.\ $C\le\Lambda$; equality forces $Q\cdot P=0$ and $r=L^2$, i.e.\ constant radius, and conversely.
\end{proof}

\begin{proposition}[Planar Kepler problem, mass $m$, strength $\kappa$]\label{bpP:p7:kepler}
Let $\mathcal D^m$ and $\mathsf Q_{m,\kappa}$ be as in Definition~\ref{bpP:p7:kepler-data}.
\begin{enumerate}[label=(\alph*)]
\item \textup{(Delaunay chart)} $\mathsf Q_{m,\kappa}$ is real analytic, $2\pi$-periodic in $\ell$ and in $g$, induces a real-analytic diffeomorphism $(\R/2\pi\Z)^2\times\{\Lambda>C>0\}\to\mathcal K^*_{m,\kappa}:=\{(q,p)\in\mathcal K_{m,\kappa}:C<\Lambda\}$ (the non-circular part), and
$\mathsf Q_{m,\kappa}^*\Bigl(\sum dp_i\wedge dq_i\Bigr)=d\Lambda\wedge d\ell+dC\wedge dg$.
\item \textup{(Identities)} For $(q,p)=\mathsf Q_{m,\kappa}(\ell,g,\Lambda,C)$, $a=\Lambda^2/(m^2\kappa)$, $e=\sqrt{1-C^2/\Lambda^2}$, $u=u_\R(\ell,e)$ \textup{(\Iref{bpI:p7:kepler-eq}(c))}:
$\ q=a\bigl((\cos u-e)e_g+\sqrt{1-e^2}\,\sin u\,J_0e_g\bigr)$, $\ \pl q=a(1-e\cos u)$, $\ \operatorname{cr}(q,p)=C$, $\ E_{m,\kappa}(q,p)=-\dfrac{\kappa^2m^3}{2\Lambda^2}$, and the elements of Definition~\ref{bpP:p7:kepler-data} are $\Lambda,C,a,e$.
\item \textup{(Mean values; includes the circular case $e=0$)} For $\Lambda>0$, $0<C\le\Lambda$ and $g\in\R$, with $(q,p)=\mathsf Q_{m,\kappa}(\ell,g,\Lambda,C)$ \textup{(extended to $C=\Lambda$ by the same formulas, \Iref{bpI:p7:circular})}:
\[\frac1{2\pi}\int_0^{2\pi}\frac{d\ell}{\pl q}=\frac1a,\qquad\frac1{2\pi}\int_0^{2\pi}\pl p^2d\ell=\frac{m^2\kappa}a,\qquad\frac1{2\pi}\int_0^{2\pi}\pl q\,d\ell=a\Bigl(1+\frac{e^2}2\Bigr).\]
\item \textup{(Hamiltonian)} $H_{\rm Kep,1}:=E_{m,\kappa}$ satisfies $H_{\rm Kep,1}\circ\mathsf Q_{m,\kappa}=h_m(\Lambda):=-\kappa^2m^3/(2\Lambda^2)$, and $h_m'(\Lambda)=\kappa^2m^3/\Lambda^3=\sqrt\kappa\,a^{-3/2}$ \textup{(Kepler's third law)}; the Hamilton equations of $h_m(\Lambda)$ in the variables $(\ell,g;\Lambda,C)$ are $\dot\ell=h_m'(\Lambda)$, $\dot g=\dot\Lambda=\dot C=0$.
\end{enumerate}
\source{Fej04r Thm.~62 and (32)--(33), p.~48 (energy $-\mu^3M^2/(2\Lambda^2)$, mean motion $\nu=M/a^{3/2}$-type formula); Chi05 \S2 for the Delaunay chart (via track~I); the dictionary is \Iref{bpI:p7:scaling}; the mean values are standard \textup{(track~I states them for $\kappa=m=1$)}.}
\end{proposition}
\begin{proof}
Throughout, $L=\Lambda/m$, $G=C/m$, $(Q,P)=\mathsf Q^{\rm I}(\ell,g,L,G)$, $(q,p)=\mathsf S_{m,\kappa}(Q,P)$; note $\Lambda>C>0$ iff $L>G>0$, i.e.\ $(L,G)\in\mathcal C_0$ of \Iref{bpI:p7:orbit-maps}.

\proofstep{(a) Chart}\label{stepP:p7:ke-chart} $\mathsf S_{m,\kappa}$ is linear and invertible, and $(\Lambda,C)\mapsto(L,G)=(\Lambda,C)/m$ is a linear isomorphism of $\{\Lambda>C>0\}$ onto $\mathcal C_0$. By \Iref{bpI:p7:chart}, $\mathsf Q^{\rm I}$ induces a real-analytic diffeomorphism of $(\R/2\pi\Z)^2\times\mathcal C_0$ onto $\mathcal C=\{E_1<0,\ \operatorname{cr}>0,\ \operatorname{cr}<L\}$ (\Iref{bpI:p7:elements}). Under $(q,p)=\mathsf S_{m,\kappa}(Q,P)$ one has $E_{m,\kappa}(q,p)=m\kappa^2E_1(Q,P)$, $\operatorname{cr}(q,p)=m\operatorname{cr}(Q,P)$ and (Step~\ref{stepP:p7:rg-scale}) $\Lambda=mL(Q,P)$; so $\mathsf S_{m,\kappa}$ maps $\mathcal C$ bijectively onto $\{E_{m,\kappa}<0,\ C>0,\ C<\Lambda\}=\mathcal K^*_{m,\kappa}$.
\proofstep{(a) The form}\label{stepP:p7:ke-form} $\sum dp_i\wedge dq_i=\sum(m\kappa\,dP_i)\wedge(dQ_i/\kappa)=-m\sum dQ_i\wedge dP_i$. By \Iref{bpI:p7:symplectic}, pulling back by $\mathsf Q^{\rm I}$ gives $-m(d\ell\wedge dL+dg\wedge dG)=m\,dL\wedge d\ell+m\,dG\wedge dg=d\Lambda\wedge d\ell+dC\wedge dg$.
\proofstep{(b)}\label{stepP:p7:ke-b} \Iref{bpI:p7:orbit-identities} for $(Q,P)$: $\pl Q=L^2(1-e\cos u)$, $E_1(Q,P)=-1/(2L^2)$, $\operatorname{cr}(Q,P)=G$, with $e=\sqrt{1-G^2/L^2}=\sqrt{1-C^2/\Lambda^2}$. With $q=Q/\kappa$: $\pl q=a(1-e\cos u)$ for $a=L^2/\kappa=\Lambda^2/(m^2\kappa)$ (this is the $a$ of Definition~\ref{bpP:p7:kepler-data}: $a=-\kappa m/(2E_{m,\kappa})$ with $E_{m,\kappa}=m\kappa^2E_1=-m\kappa^2/(2L^2)$ gives $a=L^2/\kappa$). The formula for $q$ is $Q/\kappa$ with $Q=\mathsf R(g)(X,Y)$, $X=L^2(\cos u-e)$, $Y=L^2\sqrt{1-e^2}\sin u$ and $\mathsf R(g)(x,y)=xe_g+yJ_0e_g$. $\operatorname{cr}(q,p)=m\operatorname{cr}(Q,P)=mG=C$; $E_{m,\kappa}(q,p)=-m\kappa^2/(2L^2)=-\kappa^2m^3/(2\Lambda^2)$.
\proofstep{(c)}\label{stepP:p7:ke-c} For $C<\Lambda$ and $C=\Lambda$ (circular), \Iref{bpI:p7:means} and \Iref{bpI:p7:circular} give $\langle1/\pl Q\rangle=1/L^2$, $\langle\pl P^2\rangle=1/L^2$, $\langle\pl Q\rangle=L^2(1+e^2/2)$. Then $1/\pl q=\kappa/\pl Q$ has mean $\kappa/L^2=1/a$; $\pl p^2=m^2\kappa^2\pl P^2$ has mean $m^2\kappa^2/L^2=m^2\kappa/a$; $\pl q=\pl Q/\kappa$ has mean $a(1+e^2/2)$.
\proofstep{(d)}\label{stepP:p7:ke-d} The value of $E_{m,\kappa}\circ\mathsf Q_{m,\kappa}$ is (b). $h'_m=\kappa^2m^3/\Lambda^3$, and $\Lambda^3=m^3\kappa^{3/2}a^{3/2}$ gives $\sqrt\kappa\,a^{-3/2}$. By (a) the chart $\mathsf Q_{m,\kappa}$ is symplectic with factor one in the variables $(\ell,g;\Lambda,C)$ (positions $(\ell,g)$, momenta $(\Lambda,C)$, $\Omega=d\Lambda\wedge d\ell+dC\wedge dg$), so the Hamilton equations of $h_m(\Lambda)$ in these variables are $\dot\ell=\partial_\Lambda h_m$, $\dot g=\partial_Ch_m=0$, $\dot\Lambda=-\partial_\ell h_m=0$, $\dot C=-\partial_gh_m=0$, and Lemma~\ref{bpP:p1:chart-field} transports them to the Kepler flow.
\end{proof}

\subsection{Kepler's equation in the eccentric longitude}\label{bpP:p7:sec-regular}

The Delaunay map is singular at $e=0$ ($g$ is undefined). The D25 formulas and the Poincar\'e chart use instead the \emph{eccentric longitude} $F=u+g$ and the \emph{eccentricity vector} $v=e\,e_g$.
Substituting $\ell=\lambda-g$, $u=F-g$ in $u-e\sin u=\ell$ gives
\[F-e\sin(F-g)=\lambda,\qquad e\sin(F-g)=v_0\sin F-v_1\cos F\quad(v=e(\cos g,\sin g)),\]
an equation \emph{regular in $v\in\R^2$}: $F-v_0\sin F+v_1\cos F=\lambda$ (\Lean{keplerMeanLongitude}).

\begin{proposition}[Regular Kepler equation]\label{bpP:p7:regular-kepler}
For $v\in\R^2$, $F\in\R$ put $\Phi(F,v)=F-v_0\sin F+v_1\cos F$ (so $\Phi$ is \Lean{keplerMeanLongitude} with $v=$ \Lean{secularEccVector}) and $J(F,v)=\partial_F\Phi=1-v_0\cos F-v_1\sin F=1-v\cdot(\cos F,\sin F)$ (\Lean{longitudeJacobian}).
\begin{enumerate}[label=(\alph*)]
\item \textup{(real)} For $\pl v<1$: $J\ge1-\pl v>0$; $F\mapsto\Phi(F,v)$ is a real-analytic increasing bijection of $\R$ with $\Phi(F+2\pi,v)=\Phi(F,v)+2\pi$. So for each $(\lambda,v)\in\R\times\{\pl v<1\}$ there is exactly one real $F=F(\lambda,v)$ with $\Phi(F,v)=\lambda$; $F$ is real analytic on $\R\times\{\pl v<1\}$, $F(\lambda+2\pi,v)=F(\lambda,v)+2\pi$, $F(\lambda,0)=\lambda$, and
\[\partial_\lambda F=\frac1J,\qquad\partial_{v_0}F=\frac{\sin F}J,\qquad\partial_{v_1}F=-\frac{\cos F}J\qquad(J=J(F(\lambda,v),v)).\]
\item \textup{(complex)} Let $\mathcal N_0=\{(\lambda,v)\in\C\times\C^2:\ \abs{\operatorname{Im}\lambda}<\tfrac14,\ \abs{v_0}+\abs{v_1}<\tfrac14\}$ and $\mathcal S=\{F\in\C:\abs{\operatorname{Im}F}\le1\}$. There is a unique holomorphic $F:\mathcal N_0\to\mathcal S$ with $\Phi(F(\lambda,v),v)=\lambda$ (with $\Phi,J$ extended to complex arguments by the same formulas); $\abs{J}\ge1-\frac14\cosh1>\frac35$ on $\mathcal N_0$, and the formulas of (a) hold. At real points $F$ is real and equals the $F$ of (a).
\item \textup{(relation to track~I)} For $v=e\,e_g$ with $e\in(0,1)$, $g\in\R$: $F(\lambda,v)=g+u_\R(\lambda-g,e)$, where $u_\R$ is the solution of Kepler's equation of \Iref{bpI:p7:kepler-eq}(c). (The composition of \Iref{bpI:p7:kepler-eq} with $v\mapsto(g,e)$ is not defined at $v=0$, which is why (b) is proved directly: it is the holomorphic continuation of this branch through $v=0$.)
\end{enumerate}
\source{\Iref{bpI:p7:kepler-eq} (Kepler's equation near $e=0$); the regular form in $v$ is the D25 definition \Lean{keplerMeanLongitude} (derived from LP17 (10)--(11)); the contraction proof is ours.}
\end{proposition}
\begin{proof}
\proofstep{(a) Monotonicity}\label{stepP:p7:rk-mono} $\abs{v\cdot(\cos F,\sin F)}\le\pl v$ (Cauchy--Schwarz), so $J\ge1-\pl v>0$. Hence $\Phi(\cdot,v)$ is strictly increasing with $\Phi(F,v)-F$ bounded ($\le\pl v$), so $\Phi(F,v)\to\pm\infty$ as $F\to\pm\infty$: a bijection of $\R$. $\Phi(F+2\pi,v)=\Phi(F,v)+2\pi$ by periodicity of $\sin,\cos$; $\Phi(F,0)=F$.
\proofstep{(a) Analyticity}\label{stepP:p7:rk-an} $\Phi$ is real analytic on $\R\times\R^2$ with $\partial_F\Phi=J\ne0$ on $\{\pl v<1\}$; \ref{bpP:p7:e-ift} gives a local real-analytic solution near any $(\lambda_0,v_0)$, which coincides with the unique global solution; so $F$ is real analytic. Differentiating $\Phi(F(\lambda,v),v)=\lambda$: $J\,\partial_\lambda F=1$; $J\,\partial_{v_0}F-\sin F=0$; $J\,\partial_{v_1}F+\cos F=0$.
\proofstep{(b) Contraction}\label{stepP:p7:rk-contr} Let $\beta_{\lambda,v}(F)=\lambda+v_0\sin F-v_1\cos F$, so $\Phi(F,v)=\lambda\iff\beta_{\lambda,v}(F)=F$. For $F=x+iy\in\mathcal S$: $\abs{\sin F}^2=\sin^2x+\sinh^2y$, $\abs{\cos F}^2=\cos^2x+\sinh^2y$, both $\le\cosh^2y\le\cosh^21$ (\Iref{stepI:p7:ke-trig}). For $(\lambda,v)\in\mathcal N_0$ then $\abs{\operatorname{Im}\beta}\le\abs{\operatorname{Im}\lambda}+(\abs{v_0}+\abs{v_1})\cosh1<\frac14+\frac14\cdot1.55<1$, so $\beta_{\lambda,v}(\mathcal S)\subset\mathcal S$, and $\abs{\beta'}=\abs{v_0\cos F-v_1\sin F}\le\frac14\cosh1<0.39$, so $\beta_{\lambda,v}$ is a $0.39$-contraction of the closed convex set $\mathcal S$ (\ref{bpP:p7:e-banach}). It has a unique fixed point $F(\lambda,v)\in\mathcal S$.
\proofstep{(b) Holomorphy, real points}\label{stepP:p7:rk-hol} $F_0=\lambda$ (which lies in $\mathcal S$) and $F_{j+1}=\beta_{\lambda,v}(F_j)$ are holomorphic on $\mathcal N_0$ with values in $\mathcal S$ and converge uniformly (geometric rate $0.39$, Banach), so $F$ is holomorphic (\ref{bpP:p7:e-hol}). For real $(\lambda,v)\in\mathcal N_0$ the real map $\beta_{\lambda,v}$ preserves $\R\subset\mathcal S$ and is a contraction there, so its fixed point is real, equals the real solution of (a) (uniqueness there), and $F$ is real. $\abs J\ge1-(\abs{v_0}+\abs{v_1})\cosh1>1-0.39>\frac35$. The derivative formulas are the same implicit differentiation as in (a), valid for holomorphic maps.
\proofstep{(c) Relation}\label{stepP:p7:rk-rel} For $v=e\,e_g$: $v_0\sin F-v_1\cos F=e(\cos g\sin F-\sin g\cos F)=e\sin(F-g)$, so $F'=g+u_\R(\lambda-g,e)$ solves $\Phi(F',v)=F'-e\sin(F'-g)=g+u_\R-e\sin u_\R=g+(\lambda-g)=\lambda$; uniqueness in (a) gives $F'=F$.
\end{proof}

\subsection{The Poincar\'e chart}\label{bpP:p7:sec-poincare}

\begin{definition}[Poincar\'e chart of one planet]\label{bpP:p7:poincare-def}
Let $m>0$ and $\kappa>0$. The \emph{Poincar\'e domain} is $\mathcal O=\{(\lambda,\xi,\Lambda,\eta)\in\R^4:\ \Lambda>0,\ \Gamma<\Lambda\}$ with $\Gamma=\Gamma(\xi,\eta)=(\xi^2+\eta^2)/2$ (open, invariant under $\lambda\mapsto\lambda+2\pi$). On $\mathcal O$ put
\begin{align*}
k&=1-\frac\Gamma\Lambda\ \in(0,1],\qquad v=\sqrt{\frac{1-\Gamma/(2\Lambda)}{\Lambda}}\,(\xi,\eta),\qquad a=\frac{\Lambda^2}{m^2\kappa},\qquad F=F(\lambda,v)\ \ \text{(Prop.~\ref{bpP:p7:regular-kepler}(a))},\\
c&=(\cos F,\sin F),\qquad J=1-v\cdot c,\\
q&=a\Bigl(k\,c-v+\frac{v\cdot c}{1+k}\,v\Bigr),\qquad p=\frac{m^2\kappa}{\Lambda J}\Bigl(J_0c-\frac{v\cdot c}{1+k}\,J_0v\Bigr),
\end{align*}
and $\Psi_{m,\kappa}(\lambda,\xi,\Lambda,\eta)=(q,p)\in\R^2\times\R^2$. For $\kappa=1$ these are \Lean{secularEccVector}, \Lean{secularRoot} ($=k=\sqrt{1-e^2}$), \Lean{secularAxis}, \Lean{longitudeJacobian}, \Lean{ellipsePosition}, \Lean{ellipseMomentum}, and $F=$ the solution of \Lean{keplerMeanLongitude}${}=\lambda$ (Proposition~\ref{bpP:p7:d25}).
\frozen{KAM.Planetary.secularEccVector, secularRoot, secularGamma, secularAxis, ellipsePosition, ellipseMomentum, keplerMeanLongitude, longitudeJacobian (for $\kappa=1$)}
\source{Fej04r (31), Thm.~62, pp.~47--48, with the normalisation of Fej07e p.~3: $r_j=\xi_j+i\eta_j=\sqrt{2\Lambda_j}\sqrt{1-\sqrt{1-e_j^2}}\,e^{i(g_j+\theta_j)}$ (planar: $\theta_j=0$); explicit regular form: LP17 (10)--(11) via the D25 definitions.}
\end{definition}

\begin{lemma}[Elementary facts on $\mathcal O$]\label{bpP:p7:o-facts}
On $\mathcal O$: $\Gamma\ge0$, $0<k\le1$, $\pl v^2=1-k^2$ (so $\pl v<1$ and $F$ is defined, real analytic, Proposition~\ref{bpP:p7:regular-kepler}(a)); $J\ge1-\pl v>0$; the radicand of $v$ is positive; $v\cdot c=e\cos(F-g)$ and $\pl v=e$ with $e=\sqrt{1-k^2}$. For $(\xi,\eta)\ne0$ write $(\xi,\eta)=\rho\,e_g$ with $\rho=\sqrt{2\Gamma}$: then $v=e\,e_g$ and $\Gamma=\Lambda-C$ with $C=\Lambda k$.
\end{lemma}
\begin{proof}
$\Gamma<\Lambda$ and $\Lambda>0$ give $k=1-\Gamma/\Lambda>0$ and $1-\Gamma/(2\Lambda)>\frac12$. $\pl v^2=\frac{1-\Gamma/(2\Lambda)}{\Lambda}2\Gamma=\frac{2\Gamma}\Lambda\bigl(1-\frac\Gamma{2\Lambda}\bigr)=1-(1-\frac\Gamma\Lambda)^2=1-k^2$. The rest is Proposition~\ref{bpP:p7:regular-kepler}(a), and for $(\xi,\eta)\ne0$: $\pl v/\rho=\sqrt{(1-\Gamma/2\Lambda)/\Lambda}$ with $\pl v=\sqrt{1-k^2}$ the eccentricity, and the direction of $v$ is that of $(\xi,\eta)$.
\end{proof}

\begin{theorem}[Poincar\'e chart, analytic through $e=0$]\label{bpP:p7:poincare}
Let $\Psi=\Psi_{m,\kappa}:\mathcal O\to\R^2\times\R^2$ be as in Definition~\ref{bpP:p7:poincare-def}.
\begin{enumerate}[label=(\alph*)]
\item \textup{(analytic)} $\Psi$ is real analytic on $\mathcal O$, including $\xi=\eta=0$; $\Psi(\lambda+2\pi,\xi,\Lambda,\eta)=\Psi(\lambda,\xi,\Lambda,\eta)$.
\item \textup{(Delaunay on $\Gamma>0$)} Let $(\xi,\eta)=\rho\,e_g$, $\rho>0$. Then $\Psi(\lambda,\xi,\Lambda,\eta)=\mathsf Q_{m,\kappa}(\lambda-g,\ g,\ \Lambda,\ \Lambda-\Gamma)$, with the eccentric longitude $F=g+u$, $u=u_\R(\lambda-g,e)$.
\item \textup{(symplectic)} $\Omega(D\Psi\,u,D\Psi\,w)=\Omega(u,w)$ on all of $\mathcal O$, i.e.\ $\Psi^*\Omega_{\rm cart}=d\Lambda\wedge d\lambda+d\eta\wedge d\xi$ with $\Omega_{\rm cart}=\sum dp_i\wedge dq_i$. Thus $\Psi$ is a $1$-chart on $\mathcal O$ in the sense of Definition~\ref{bpP:p1:chart}, and $D\Psi$ is invertible.
\item \textup{(elements)} $\Psi(\mathcal O)\subset\mathcal K_{m,\kappa}$ and for $(q,p)=\Psi(\lambda,\xi,\Lambda,\eta)$: $E_{m,\kappa}(q,p)=-\kappa^2m^3/(2\Lambda^2)$, $a(q,p)=\Lambda^2/(m^2\kappa)$, the action of Definition~\ref{bpP:p7:kepler-data} is $\Lambda$, $\operatorname{cr}(q,p)=\Lambda-\Gamma=\Lambda k$ and, the Laplace--Runge--Lenz vector satisfies $\kappa^{-1}\bigl[\bigl(\frac{\pl p^2}{m^2}-\frac\kappa{\pl q}\bigr)q-\frac{q\cdot p}{m^2}p\bigr]=v$ (for $\kappa=1$ this is \Lean{keplerEccVector}${}=v$).
\item \textup{(bijectivity modulo $2\pi$)} $\Psi:\mathcal O\to\mathcal K_{m,\kappa}$ is surjective, and $\Psi(\lambda',\xi',\Lambda',\eta')=\Psi(\lambda,\xi,\Lambda,\eta)$ iff $(\xi',\Lambda',\eta')=(\xi,\Lambda,\eta)$ and $\lambda'-\lambda\in2\pi\Z$. The induced map $\mathcal O/(2\pi\Z\times\{0\})\to\mathcal K_{m,\kappa}$ (quotient by the $\lambda$-translations) is a real-analytic diffeomorphism. Its inverse is $\Lambda=m\sqrt{\kappa a}$, $\Gamma=\Lambda-\operatorname{cr}(q,p)$, $v=$ the rescaled LRL vector of (d), $(\xi,\eta)=v/\sqrt{(1-\Gamma/(2\Lambda))/\Lambda}$, $c=\bigl(k\,\mathrm{Id}+\tfrac1{1+k}vv^{\top}\bigr)^{-1}(q/a+v)$, $F=\arg c$, $\lambda=\Phi(F,v)$.
\item \textup{(circular torus)} At $\xi=\eta=0$: $v=0$, $k=1$, $F=\lambda$, $q=a(\cos\lambda,\sin\lambda)$, $p=m\sqrt{\kappa/a}\,(-\sin\lambda,\cos\lambda)$, $a=\Lambda^2/(m^2\kappa)$: for $\Lambda=m\sqrt{\kappa a}$ this is \Lean{circularPoint} (with $\phi=\lambda/2\pi$, one planet).
\end{enumerate}
\end{theorem}
\begin{proof}
\proofstep{(a)}\label{stepP:p7:po-an} $k$, $v$ (positive radicand, Lemma~\ref{bpP:p7:o-facts}) and $J$ are real analytic, $F$ is real analytic (Proposition~\ref{bpP:p7:regular-kepler}(a) composed with the analytic $v(\xi,\eta,\Lambda)$, $\pl v<1$), $1+k\ge1$, $\Lambda J>0$; $q,p$ are rational expressions in these and in $\cos F,\sin F$. Periodicity: $F(\lambda+2\pi,v)=F+2\pi$ and $c$ is $2\pi$-periodic in $F$.
\proofstep{(b) Delaunay on $\Gamma>0$}\label{stepP:p7:po-del} Put $\ell=\lambda-g$, $C=\Lambda-\Gamma=\Lambda k$, $L=\Lambda/m$, $G=C/m=Lk$. Then $\Lambda>C>0$, $e=\sqrt{1-C^2/\Lambda^2}=\sqrt{1-k^2}$ and $v=e\,e_g$ (Lemma~\ref{bpP:p7:o-facts}), so by Proposition~\ref{bpP:p7:regular-kepler}(c) $F=g+u$, $u=u_\R(\ell,e)$. Write $\hat e=e_g$, $\hat e^\perp=J_0e_g$ and $c=\cos u\,\hat e+\sin u\,\hat e^\perp$ ($F=u+g$). Then $v\cdot c=e\cos u$, $J=1-e\cos u$ and, using $e^2/(1+k)=(1-k^2)/(1+k)=1-k$,
\[kc-v+\frac{v\cdot c}{1+k}v=(k\cos u-e+(1-k)\cos u)\hat e+k\sin u\,\hat e^\perp=(\cos u-e)\hat e+k\sin u\,\hat e^\perp,\]
so $q=a((\cos u-e)\hat e+k\sin u\,\hat e^\perp)$, which is the position of Proposition~\ref{bpP:p7:kepler}(b). Likewise $J_0c=\cos u\,\hat e^\perp-\sin u\,\hat e$, $J_0v=e\hat e^\perp$ and
\[J_0c-\frac{v\cdot c}{1+k}J_0v=-\sin u\,\hat e+\bigl(1-\tfrac{e^2}{1+k}\bigr)\cos u\,\hat e^\perp=-\sin u\,\hat e+k\cos u\,\hat e^\perp=\mathsf R(g)\binom{-\sin u}{k\cos u}.\]
By \Iref{bpI:p7:orbit-maps}, $P=\mathsf R(g)\binom{-\sin u}{k\cos u}/(LJ)$ and $p=m\kappa P=\frac{m\kappa}{LJ}(\cdots)=\frac{m^2\kappa}{\Lambda J}(\cdots)$, which is the $p$ of the definition. Hence $\Psi=\mathsf S_{m,\kappa}\circ\mathsf Q^{\rm I}(\ell,g,L,G)=\mathsf Q_{m,\kappa}(\ell,g,\Lambda,C)$.
\proofstep{(c) Symplecticity on $\{\Gamma>0\}$}\label{stepP:p7:po-symp} On the open set $\mathcal A=\{(\xi,\eta)\ne0\}\subset\mathcal O$ choose locally a real-analytic branch $g(\xi,\eta)=\arg(\xi+i\eta)$ and put $\rho=\sqrt{2\Gamma}$. By (b), $\Psi=\mathsf Q_{m,\kappa}\circ T$ with $T(\lambda,\xi,\Lambda,\eta)=(\lambda-g,g,\Lambda,\Lambda-\Gamma)$ (real analytic, values in $\mathcal D^m$). By Proposition~\ref{bpP:p7:kepler}(a), $\Psi^*\Omega_{\rm cart}=T^*(d\Lambda\wedge d\ell+dC\wedge dg)=d\Lambda\wedge(d\lambda-dg)+(d\Lambda-d\Gamma)\wedge dg=d\Lambda\wedge d\lambda-d\Gamma\wedge dg$. Since $\xi=\rho\cos g$, $\eta=\rho\sin g$: $d\xi\wedge d\eta=\rho\,d\rho\wedge dg=d\Gamma\wedge dg$, i.e.\ $d\eta\wedge d\xi=-d\Gamma\wedge dg$. So $\Psi^*\Omega_{\rm cart}=d\Lambda\wedge d\lambda+d\eta\wedge d\xi$ on $\mathcal A$.
\proofstep{(c) All of $\mathcal O$}\label{stepP:p7:po-dense} The matrix identity $D\Psi^{\top}M_0D\Psi=M_0$ (equivalent to the form identity, \ref{bpP:p7:e-forms}) is an identity between continuous functions on $\mathcal O$ ($\Psi$ is $C^1$ there by (a)), valid on the dense subset $\mathcal A$ (the complement $\{\xi=\eta=0\}$ has empty interior); \ref{bpP:p7:e-dense}. This is the chart property of Definition~\ref{bpP:p1:chart} with $c=1$ (the factor is $1$ because the source form is the frozen $\Omega$ in the order $(\lambda,\xi;\Lambda,\eta)$). Invertibility: Lemma~\ref{bpP:p1:omega-algebra}(b).
\proofstep{(d) Elements}\label{stepP:p7:po-el} On $\mathcal A$ by (b) and Proposition~\ref{bpP:p7:kepler}(b): $E_{m,\kappa}=-\kappa^2m^3/(2\Lambda^2)<0$, $\operatorname{cr}=\Lambda k>0$, $q\ne0$ (as $\pl q=a(1-e\cos u)>0$), hence $\Psi(\mathcal A)\subset\mathcal K_{m,\kappa}$, and $a(q,p)=-\kappa m/(2E)=\Lambda^2/(m^2\kappa)$, action $m\sqrt{\kappa a}=\Lambda$. At $\xi=\eta=0$ the same identities are \Iref{bpI:p7:circular} (with $e=0$, $k=1$, $u=F=\lambda$) through the scaling $\mathsf S_{m,\kappa}$ of Step~\ref{stepP:p7:ke-b}; they can also be checked directly from (f). The LRL identity: for $(Q,P)=\mathsf S_{m,\kappa}^{-1}(q,p)$ one computes $\kappa^{-1}\bigl[(\frac{\pl p^2}{m^2}-\frac\kappa{\pl q})q-\frac{q\cdot p}{m^2}p\bigr]=(\pl P^2-\frac1{\pl Q})Q-(Q\cdot P)P=:A$. On $\mathcal A$, in the orbit frame $(\hat e,\hat e^\perp)$ of Step~\ref{stepP:p7:po-del} with $c_u=\cos u$, $s_u=\sin u$, $J=1-ec_u$, $L=\Lambda/m$ (\Iref{bpI:p7:orbit-identities}): $\pl P^2=\frac{1+ec_u}{L^2J}$, $\frac1{\pl Q}=\frac1{L^2J}$, $Q=L^2(c_u-e,\,ks_u)$, $Q\cdot P=Les_u$, $P=\frac1{LJ}(-s_u,kc_u)$, hence
$A=\frac{ec_u}J(c_u-e,\,ks_u)-\frac{es_u}{J}(-s_u,\,kc_u)=\bigl(\frac eJ(c_u^2-ec_u+s_u^2),\,0\bigr)=(e,0)$ in the frame, i.e.\ $A=e\,\hat e=v$. At $\xi=\eta=0$ both sides vanish ($\pl Q=L^2$, $\pl P^2=1/L^2$, $Q\cdot P=0$ by (f)).
\proofstep{(e) Bijectivity}\label{stepP:p7:po-bij} \emph{Injectivity modulo $2\pi$ in $\lambda$.} Let $\Psi(w')=\Psi(w)$, $w,w'\in\mathcal O$. By (d) $\Lambda'=\Lambda$ ($\Lambda=m\sqrt{\kappa a}$ is a function of $(q,p)$) and $\Gamma'=\Lambda-\operatorname{cr}=\Gamma$. If $\Gamma=0$: $\xi=\eta=\xi'=\eta'=0$ and by (f) $\Psi=(a(\cos\lambda,\sin\lambda),\cdots)$ so $\lambda'-\lambda\in2\pi\Z$. If $\Gamma>0$: $w,w'\in\mathcal A$, $(\ell,g)$, $(\ell',g')$ as in (b), and by \Iref{bpI:p7:chart} (injectivity modulo $2\pi$ of $\mathsf Q^{\rm I}$, transported by Proposition~\ref{bpP:p7:kepler}(a)) $g'-g\in2\pi\Z$, $\ell'-\ell\in2\pi\Z$. Then $(\xi',\eta')=\rho\,e_{g'}=\rho\,e_g=(\xi,\eta)$ and $\lambda'=\ell'+g'\in\lambda+2\pi\Z$. The converse is (a).
\emph{Surjectivity.} Let $(q,p)\in\mathcal K_{m,\kappa}$, $\Lambda$, $C$ its elements, $\Gamma=\Lambda-C\in[0,\Lambda)$ (Lemma~\ref{bpP:p7:range}). If $C<\Lambda$: $(q,p)\in\mathcal K^*_{m,\kappa}$ is $\mathsf Q_{m,\kappa}(\ell,g,\Lambda,C)$ by Proposition~\ref{bpP:p7:kepler}(a), and with $(\xi,\eta)=\sqrt{2\Gamma}\,e_g$, $\lambda=\ell+g$, (b) gives $\Psi(\lambda,\xi,\Lambda,\eta)=(q,p)$. If $C=\Lambda$ (so $\Gamma=0$): in the scaling of Step~\ref{stepP:p7:rg-scale}, $(Q,P)$ has $\operatorname{cr}=L$ and the equality case of Step~\ref{stepP:p7:rg-lagrange} gives $\pl Q=L^2$, $Q\cdot P=0$, $\pl P=\operatorname{cr}/\pl Q=1/L$; so $Q=L^2(\cos\lambda,\sin\lambda)$ for some $\lambda$ and $P=\frac1L(-\sin\lambda,\cos\lambda)$ (orthogonal to $Q$, of length $1/L$, positive cross product $\operatorname{cr}(Q,P)=L>0$). That is $(q,p)=\Psi(\lambda,0,\Lambda,0)$ by (f).
\emph{Diffeomorphism.} $D\Psi$ is invertible (c); $\Psi$ is a bijective, real-analytic local diffeomorphism from the quotient by the lattice $2\pi\Z$ in $\lambda$ onto the open set $\mathcal K_{m,\kappa}$ (\ref{bpP:p7:e-ift}). Inverse formulas: $\Lambda,\Gamma$ as in (d); $v$ by the LRL identity (d); $(\xi,\eta)=v/s$ with $s=\sqrt{(1-\Gamma/(2\Lambda))/\Lambda}>0$ (defined also at $v=0$); $\frac qa=kc-v+\frac{v\cdot c}{1+k}v=(k\,\mathrm{Id}+\frac1{1+k}vv^{\top})c-v$ and $k\,\mathrm{Id}+\frac1{1+k}vv^{\top}$ has eigenvalues $k$ and $k+\pl v^2/(1+k)=1$, hence is invertible; $c$ is then the unit vector $\cos F,\sin F$, and $\lambda=\Phi(F,v)$ modulo $2\pi$.
\proofstep{(f) Circular points}\label{stepP:p7:po-circ} $\Gamma=0$ gives $v=0$, $k=1$, $F=\lambda$ (Proposition~\ref{bpP:p7:regular-kepler}(a)), $J=1$, $q=a(c-0+0)=a(\cos\lambda,\sin\lambda)$, $p=\frac{m^2\kappa}\Lambda J_0c=m\sqrt{\kappa/a}\,J_0c$ since $\Lambda=m\sqrt{\kappa a}$. Compare Definition~\ref{bpP:p0:circular}.
\end{proof}

\begin{proposition}[Holomorphic extension near the circular orbits]\label{bpP:p7:holomorphic}
Let $\Lambda^\circ>0$ and
\[\mathcal N(\Lambda^\circ)=\Bigl\{(\lambda,\xi,\Lambda,\eta)\in\C^4:\ \abs{\operatorname{Im}\lambda}<\tfrac14,\ \abs{\Lambda-\Lambda^\circ}<\tfrac12\Lambda^\circ,\ \abs\xi,\abs\eta<\tfrac1{16}\sqrt{\Lambda^\circ}\Bigr\}.\]
The formulas of Definition~\ref{bpP:p7:poincare-def} (with $\Gamma=(\xi^2+\eta^2)/2$, $F=F(\lambda,v)$ of Proposition~\ref{bpP:p7:regular-kepler}(b)) define a holomorphic map $\Psi:\mathcal N(\Lambda^\circ)\to\C^4$, $2\pi$-periodic in $\lambda$, real at real points, with $\mathcal N(\Lambda^\circ)\cap\R^4\subset\mathcal O$ and $D\Psi^{\top}M_0D\Psi=M_0$ at real points. Hence it satisfies the hypotheses of Lemma~\ref{bpP:p1:hol-chart} with $c=1$.
\end{proposition}
\begin{proof}
\proofstep{Sizes}\label{stepP:p7:ho-size} On $\mathcal N(\Lambda^\circ)$: $\abs\Lambda>\frac12\Lambda^\circ$ and $\abs{\arg\Lambda}\le\arcsin\frac12=\frac\pi6$; $\abs\Gamma<\frac1{256}\Lambda^\circ$, so $\abs{\Gamma/\Lambda}<\frac1{128}$, hence $\abs{\arg(1-\Gamma/2\Lambda)}<0.01$, $\abs{1-\Gamma/2\Lambda}<1.01$, $\operatorname{Re}k>0.99$ ($k=1-\Gamma/\Lambda$), $\abs{1+k}>1.9$. The radicand $(1-\Gamma/2\Lambda)/\Lambda$ has argument of modulus $<\frac\pi6+0.01<\frac\pi2$, so its principal square root $s$ is holomorphic (\ref{bpP:p7:e-hol}), and $\abs s\le\sqrt{1.01\cdot2/\Lambda^\circ}<1.5/\sqrt{\Lambda^\circ}$. Therefore $\abs{v_0}+\abs{v_1}\le\abs s(\abs\xi+\abs\eta)<1.5\cdot\frac2{16}<\frac14$, $(\lambda,v)\in\mathcal N_0$, and $F$ is holomorphic with $\abs J>\frac35$ (Proposition~\ref{bpP:p7:regular-kepler}(b)).
\proofstep{Holomorphy, reality, periodicity}\label{stepP:p7:ho-hol} $q,p$ are rational in $(\Lambda,k,v,\cos F,\sin F,J^{-1},(1+k)^{-1})$, all holomorphic with nonvanishing denominators; periodicity as in Step~\ref{stepP:p7:po-an}; at real points everything is real and $\Gamma<\Lambda$, so the point lies in $\mathcal O$ and the real maps agree. The matrix identity at real points is Theorem~\ref{bpP:p7:poincare}(c) (the complex derivative restricts to the real one, \Kref{bp:p1:real-restriction}).
\end{proof}
\begin{remark}\label{bpP:p7:ho-rem}
The neighbourhood is small in $v$ by design (a contraction in $\abs{v_0}+\abs{v_1}<\frac14$); this is all that the circular-orbit analysis (Parts~2, 5, 8) needs: $\abs{\xi},\abs\eta$ of size $\ll\sqrt\Lambda$ is eccentricity $e\lesssim\frac14$ (cf.\ $\pl v=e$, Lemma~\ref{bpP:p7:o-facts}). For larger eccentricities (not needed in this track) the complex tube of \Iref{bpI:p7:kepler-eq} covers $e\le e_1<1$ in the Delaunay form away from $e=0$.
\end{remark}

\begin{remark}[Normalisation, sign conventions, what the source does not prove]\label{bpP:p7:normalisation}
(1) \emph{Normalisation.} Fej04r (31) as printed has $r_j=\xi_j+i\eta_j=\sqrt{\Lambda_j}\sqrt{1-\sqrt{1-\varepsilon_j^2}}\,e^{i(g_j+\theta_j)}$ with only $\sqrt{\Lambda_j}$, for which the chart is merely conformally symplectic (factor $2$ in $d\xi\wedge d\eta$). Fej07e p.~3 (``sans le facteur $\sqrt2$, elles ne sont que conform\'ement symplectiques'') corrects it to $\sqrt{2\Lambda_j}$, i.e.\ $\frac12(\xi^2+\eta^2)=\Lambda-C=\Gamma$ (Lemma~\ref{bpP:p7:o-facts}), which is exactly \Lean{secularGamma}: $\Gamma=R(\eta,\xi)$ with Definition~\ref{bpP:p0:oscaction}. With this normalisation the oscillator frequencies of the secular quadratic form are the diagonal values of $\partial^2\langle H\rangle/\partial\Gamma^2$ with no extra factor, which is what later fixes the frozen $\Omega=2\sigma$.
(2) \emph{Sign.} Fej04r Thm.~62 writes $\omega=\sum d\lambda_j\wedge d\Lambda_j+d\xi_j\wedge d\eta_j$ (position$\wedge$momentum), and Fej07e p.~3 warns that signs vary (``conventions de signe variables quant \`a la forme symplectique'') and gives the a-posteriori test (the perihelion of the larger planet precesses forward). The frozen order $(\lambda,\xi;\Lambda,\eta)$ makes $\xi$ the position and $\eta$ the momentum, $\Omega=\sum dz_2\wedge dz_1$; with that order the identity of Theorem~\ref{bpP:p7:poincare}(c) forces $(\xi,\eta)=\rho(\cos g,\sin g)$ (the argument of $\xi+i\eta$ is $+g$, as in Fej07e). If one writes the same chart with the roles of $\xi,\eta$ exchanged ($\eta$ position), the formula becomes $(\xi,\eta)=\rho(\cos g,-\sin g)$; this is only a relabelling and is not used.
(3) \emph{Source gap.} Fej04r Thm.~62 gives no proof (``Voir [Poin05, Chap.~III], ou [Che88, Appendice]'', neither pinned); in particular the analyticity \emph{through the circular orbits} is the whole content of the theorem. The proof above is ours and uses only explicit formulas: the regular Kepler equation (Proposition~\ref{bpP:p7:regular-kepler}), the Delaunay chart of track~I on $\Gamma>0$, and a density argument for the symplectic identity at $\Gamma=0$.
\sourcegap{Fej04r Thm.~62 (p.~48) cites Poincar\'e and Chenciner, not pinned; analyticity at $e=0$ and bijectivity at the circular orbits are proved here from explicit formulas (Theorem~\ref{bpP:p7:poincare}(a),(c),(e)).}
\end{remark}

\subsection{The product chart for the two physical models}\label{bpP:p7:sec-product}

Fix $n\ge1$, $G,m_0>0$ ($\kappa=Gm_0$) and $m=(m_i)\in(0,\infty)^n$. Coordinates $w=(w_1,w_2)\in\Ph{2n}$ with $w_1=(\lambda,\xi)$ and $w_2=(\Lambda,\eta)$, $\lambda,\xi,\Lambda,\eta\in\R^n$ (fast block first: the order of \Lean{assemble}, Definition~\ref{bpP:p0:blocks}, with $a=b=n$). Physical points are $(q,p)\in\Ph{2n}$ with $q^{(i)}=\operatorname{body}(q,i)$.

\begin{definition}[Product chart]\label{bpP:p7:product-def}
$\mathcal O_n=\{w:\ (\lambda_i,\xi_i,\Lambda_i,\eta_i)\in\mathcal O\ \forall i\}=\{\Lambda_i>0,\ \tfrac12(\xi_i^2+\eta_i^2)<\Lambda_i\ \forall i\}$ and $\Psi^{(n)}(w)=(q,p)$ with $(q^{(i)},p^{(i)})=\Psi_{m_i,\kappa}(\lambda_i,\xi_i,\Lambda_i,\eta_i)$. $h(\Lambda)=-\sum_i\kappa^2m_i^3/(2\Lambda_i^2)$ and $\mathcal K^{(n)}=\{(q,p):\ (q^{(i)},p^{(i)})\in\mathcal K_{m_i,\kappa}\ \forall i\}$; for $\kappa=1$, $\mathcal K^{(n)}=\{\operatorname{ProgradeElliptic}\ m\ z\}$ (\Lean{ProgradeElliptic}: $0<m_i$, $\pl{q^{(i)}}>0$, \Lean{keplerEnergy}${}<0$, \Lean{planeCross}${}>0$).
\end{definition}

\begin{theorem}[Product Poincar\'e chart]\label{bpP:p7:product-chart}
\begin{enumerate}[label=(\alph*)]
\item $\Psi^{(n)}$ is a real-analytic $1$-chart on the open set $\mathcal O_n$ (Definition~\ref{bpP:p1:chart}, $c=1$), $\mathcal O_n+((2\pi k,0),0)=\mathcal O_n$, $\Psi^{(n)}(w+((2\pi k,0),0))=\Psi^{(n)}(w)$ for $k\in\Z^n$, and $\Psi^{(n)}(w')=\Psi^{(n)}(w)$ iff $w'=w+((2\pi k,0),0)$ for some $k\in\Z^n$. It maps $\mathcal O_n$ onto $\mathcal K^{(n)}$, and induces a real-analytic diffeomorphism of $\mathcal O_n/\{(2\pi k,0,0,0)\}$ onto $\mathcal K^{(n)}$.
\item \textup{(Decoupling)} $H_{\rm Kep}\circ\Psi^{(n)}=h(\Lambda)$ on $\mathcal O_n$, with $H_{\rm Kep}$ of Lemma~\ref{bpP:p6:split}; $h$ does not depend on $\lambda,\xi,\eta$, nor on any coupling parameter. $\partial h/\partial\Lambda_i=\nu_i(\Lambda_i):=\kappa^2m_i^3/\Lambda_i^3$ and $\partial^2h/\partial\Lambda_i^2=-3\kappa^2m_i^3/\Lambda_i^4\ne0$; $h$ is real analytic on $\{\Lambda_i>0\}$.
\item \textup{(The two models share the chart)} Let $\Psi^{-1}(W)=\{w\in\mathcal O_n:\Psi^{(n)}(w)\in W\}$ (open, $W$ of Lemma~\ref{bpP:p6:domain}) and
\[\hat f_{\rm mov}=f_{\rm mov}\circ\Psi^{(n)},\qquad\hat f_{\rm fix}=f_{\rm fix}\circ\Psi^{(n)}\qquad(f\ \text{of Lemma~\ref{bpP:p6:split}}).\]
These are real analytic on $\Psi^{-1}(W)$ and do not depend on $\mu$ or $\varepsilon$. For all $\mu,\varepsilon$ and $w\in\mathcal O_n$: $H_{\rm mov}^{\Psi}(w)=h(\Lambda)+\mu\hat f_{\rm mov}(w)$ and $H_{\rm fix}^{\Psi}(w)=h(\Lambda)+\varepsilon\hat f_{\rm fix}(w)$ (the Lean total functions; $H^\Psi=H\circ\Psi^{(n)}$ because $c=1$). On $\Psi^{-1}(W)$ both are differentiable and Lemma~\ref{bpP:p1:chart-field} gives $D\Psi^{(n)}X_{H^{\Psi}}=X_H\circ\Psi^{(n)}$; orbits transfer by Corollary~\ref{bpP:p1:chart-orbit}.
\item \textup{(Hamilton equations in the chart)} On $\Psi^{-1}(W)$ the equations $\dot w=X_{H^\Psi}(w)$ read $\dot\lambda_i=\nu_i(\Lambda_i)+\partial_{\Lambda_i}(\mu\hat f)$, $\dot\Lambda_i=-\mu\partial_{\lambda_i}\hat f$, $\dot\xi_i=\mu\partial_{\eta_i}\hat f$, $\dot\eta_i=-\mu\partial_{\xi_i}\hat f$ (with $\mu\hat f$ replaced by $\varepsilon\hat f_{\rm fix}$ for $H_{\rm fix}$).
\end{enumerate}
\end{theorem}
\begin{proof}
\proofstep{Product structure}\label{stepP:p7:pc-prod} The $i$-th pair $(q^{(i)},p^{(i)})$ depends only on $(\lambda_i,\xi_i,\Lambda_i,\eta_i)$, and $\Omega(v,w)=\sum_i\Omega_i$ with $\Omega_i(v,w)=v_{2,i}\cdot w_{1,i}-v_{1,i}\cdot w_{2,i}$ over the coordinates $(\lambda_i,\xi_i;\Lambda_i,\eta_i)$ (positions $\lambda_i,\xi_i$) resp.\ $(q^{(i)};p^{(i)})$; the index bookkeeping ($k+2i$, \Lean{finProdFinEquiv}) is a permutation applied to positions and to momenta alike, which does not change $\Omega=\sum dz_2\wedge dz_1$. Hence $D\Psi^{(n)}$ is block diagonal with blocks $D\Psi_{m_i,\kappa}$ each satisfying the identity of Theorem~\ref{bpP:p7:poincare}(c); the sum gives $\Omega(D\Psi u,D\Psi w)=\Omega(u,w)$. Analyticity, lattice invariance, bijectivity, the diffeomorphism and $\mathcal O_n$ open are blockwise Theorem~\ref{bpP:p7:poincare}(a),(c),(e) (the lattice $2\pi\Z^n$ in $\lambda$ acts coordinatewise).
\proofstep{Decoupling}\label{stepP:p7:pc-dec} $H_{\rm Kep}(q,p)=\sum_i\bigl[\pl{p_i}^2/(2m_i)-\kappa m_i/\pl{q_i}\bigr]=\sum_iE_{m_i,\kappa}(q_i,p_i)$ (Remark~\ref{bpP:p6:kepler-part}), and $E_{m_i,\kappa}\circ\Psi_{m_i,\kappa}=-\kappa^2m_i^3/(2\Lambda_i^2)$ by Theorem~\ref{bpP:p7:poincare}(d). Derivatives are elementary.
\proofstep{Shared chart}\label{stepP:p7:pc-share} By Lemma~\ref{bpP:p6:split}, $H_{\rm mov}=H_{\rm Kep}+\mu f_{\rm mov}$, $H_{\rm fix}=H_{\rm Kep}+\varepsilon f_{\rm fix}$ on all of $\Ph{2n}$ (total functions), with $H_{\rm Kep}$, $f_{\rm mov}$, $f_{\rm fix}$ independent of $\mu,\varepsilon$; composing with $\Psi^{(n)}$ gives the formulas. $f_{\rm mov},f_{\rm fix}$ are real analytic on $W$ (finite sums of polynomials and $1/\pl{\cdot}$ with positive radicand, Lemma~\ref{bpP:p6:domain}, whose proof gives analyticity as well as $C^\infty$); $\Psi^{(n)}$ is analytic; the composition is analytic by \ref{bpP:p1:ext-comp}. The domain $\Psi^{-1}(W)$ is open since $W$ is open and $\Psi^{(n)}$ continuous. The field statements are Lemma~\ref{bpP:p1:chart-field}, Corollary~\ref{bpP:p1:chart-orbit}.
\proofstep{(d)}\label{stepP:p7:pc-eq} Definition~\ref{bpP:p0:hamsymp}: $X_G=(\partial_{z_2}G,-\partial_{z_1}G)$ with $z_1=(\lambda,\xi)$, $z_2=(\Lambda,\eta)$.
\end{proof}

\begin{remark}[Why the same chart serves both models]\label{bpP:p7:same-chart}
The Kepler part of $H_{\rm mov}$ is $H_{\rm Kep}$ with the \emph{planet} mass $m_i$ and centre mass $m_0$, not Fej04r's reduced mass $\mu_i=m_0m_i/(m_0+\mu m_i)$ and $M_i=m_0+\mu m_i$ (Remark~\ref{bpP:p6:kepler-part}): the term $\mu\pl{p_i}^2/(2m_0)$ is part of $\mu f_{\rm mov}$. The chart is therefore independent of $\mu$ and equal to the chart of the fixed-Sun problem; this is what makes $H_{\rm mov}=h+\mu f$ and $H_{\rm fix}=h+\varepsilon f'$ with coupling-independent $f$, $f'$. The difference $f_{\rm mov}-f_{\rm fix}=\pl P^2/(2m_0)$, $P=\sum_ip_i$, is analytic on all of $\Ph{2n}$ (a polynomial) and its $\lambda$-average is computed in Part~8 from Proposition~\ref{bpP:p7:kepler}(c) and Fej04r Lemma~64.
\end{remark}

\begin{lemma}[The tube in chart coordinates]\label{bpP:p7:tube}
Let $a_i>0$ and $\Lambda_i^*=m_i\sqrt{\kappa a_i}$. For every $\delta>0$ there is $\eta_\delta>0$ such that for all $\eta\in(0,\eta_\delta]$: the circular tube $U_\eta$ (Definition~\ref{bpP:p0:circular}, with $G$, $m_0$, $a$ as in the definition) is contained in $\mathcal K^{(n)}$, and every $w\in\mathcal O_n$ with $\Psi^{(n)}(w)\in U_\eta$ satisfies $\abs{\Lambda_i-\Lambda_i^*}<\delta$ and $\tfrac12(\xi_i^2+\eta_i^2)<\delta$ for all $i$. Moreover $\operatorname{circ}(\phi)=\Psi^{(n)}\bigl((2\pi\phi,0),(\Lambda^*,0)\bigr)$.
\end{lemma}
\begin{proof}
\proofstep{Circular points}\label{stepP:p7:tu-circ} Theorem~\ref{bpP:p7:poincare}(f) blockwise: with $\Lambda_i=\Lambda_i^*$, $a=a_i$, $\xi=\eta=0$ and $\lambda_i=2\pi\phi_i$ the image is $\operatorname{circ}(\phi)$.
\proofstep{Continuity of the elements}\label{stepP:p7:tu-cont} On the open set $\mathcal K^{(n)}$ the functions $\Lambda_i(z)=m_i\sqrt{-\kappa^2m_i/(2E_{m_i,\kappa})}$ and $\Gamma_i(z)=\Lambda_i(z)-\operatorname{cr}(q^{(i)},p^{(i)})$ are continuous, and by Theorem~\ref{bpP:p7:poincare}(d) they equal $\Lambda_i,\Gamma_i$ of any preimage. The set $V_\delta=\{z\in\mathcal K^{(n)}:\abs{\Lambda_i(z)-\Lambda_i^*}<\delta,\ \Gamma_i(z)<\delta\ \forall i\}$ is open and contains the compact torus $T=\{\operatorname{circ}(\phi)\}$ (Step~\ref{stepP:p7:tu-circ}: there $\Lambda_i=\Lambda_i^*$, $\Gamma_i=0$; $T$ is the continuous image of $[0,1]^n$). By \ref{bpP:p7:e-comp} there is $\eta_\delta>0$ with $\{z:\operatorname{dist}_\infty(z,T)<\eta_\delta\}\subset V_\delta$; this set is $U_{\eta_\delta}$ (Definition~\ref{bpP:p0:circular}). For $\eta\le\eta_\delta$, $U_\eta\subset V_\delta\subset\mathcal K^{(n)}$, and a preimage $w$ of $z\in U_\eta$ has $\Lambda_i(w)=\Lambda_i(z)$, $\tfrac12(\xi_i^2+\eta_i^2)=\Gamma_i(z)<\delta$.
\end{proof}

\begin{lemma}[Lattice compatibility for Part~1]\label{bpP:p7:lattice}
Let $\Theta_n:\{(\varphi,\psi;I,R):\ I_i>0,\ 0<R_i<I_i\}\to\mathcal O_n$ be the oscillator polar chart
\[\Theta_n((\varphi,\psi),(I,R))=\bigl((\varphi,\ \sqrt{2R}\sin\psi),\ (I,\ \sqrt{2R}\cos\psi)\bigr)\]
(componentwise square roots and sines; the formula of Remark~\ref{bpP:p1:typical}(1): $\xi=\sqrt{2R}\sin\psi$ is the position, $\eta=\sqrt{2R}\cos\psi$ the momentum; $\Gamma_i=R_i$, $\Lambda_i=I_i$), and $\Phi_n=\Psi^{(n)}\circ\Theta_n$ on $O_0=\{I_i>0,0<R_i<I_i\}$. Then:
\begin{enumerate}[label=(\alph*)]
\item $\Theta_n$ is a $1$-chart on $O_0$; $\Phi_n$ is a $1$-chart on $O_0$; $O_0+\iota'_k=O_0$ for the full lattice shift $\iota'_k=((2\pi k,0),0)$ of all $2n$ positions, $k\in\Z^{2n}$ ($\iota'_k$ adds $2\pi k$ to $(\varphi,\psi)$).
\item $\Phi_n(z+\iota'_k)=\Phi_n(z)$ for $z\in O_0$, $k\in\Z^{2n}$; and if $z,z'\in O_0$, $\Phi_n(z')=\Phi_n(z)$ then $z'=z+\iota'_k$ for some $k\in\Z^{2n}$. This is (Q1) (with $\sigma_k=0$) and (Q2) of Theorem~\ref{bpP:p1:torus-transport} in the radian convention; the period-one version is (c).
\item In the period-one torus parameter of the frozen predicates: $\Phi_n^{(1)}:=\Phi_n\circ S_{2\pi}$ (with $S_{2\pi}$ of Lemma~\ref{bpP:p1:period-convert}) is a $2\pi$-chart on $S_{2\pi}^{-1}(O_0)$ satisfying (Q1) with $\sigma_k=0$ for the lattice $\iota_k=(k,0)$, $k\in\Z^{2n}$, and (Q2); in particular Corollary~\ref{bpP:p1:phys-pipeline} applies to $\Phi_n$ and Theorem~\ref{bpP:p1:torus-transport}(b) to $\Phi_n^{(1)}$.
\item Under $\Theta_n$ the pulled-back Hamiltonians are $H^{\Phi_n}=H\circ\Psi^{(n)}\circ\Theta_n=h(I)+\mu\hat f\circ\Theta_n$ ($I=\Lambda$), i.e.\ in the variables $(\varphi,\psi;I,R)$ the unperturbed part is $h(I)$, independent of $\psi$, $R$.
\end{enumerate}
\end{lemma}
\begin{proof}
\proofstep{(a) Polar chart}\label{stepP:p7:la-polar} $\Theta_n$ is blockwise $(\psi,R)\mapsto(\xi,\eta)=\sqrt{2R}(\sin\psi,\cos\psi)$, with $d\eta\wedge d\xi=dR\wedge d\psi$ (computed: $d\eta\wedge d\xi=\rho\,d\rho\wedge d\psi=dR\wedge d\psi$ for $\rho=\sqrt{2R}$), real analytic for $R>0$, mapping into $\mathcal O_n$ since $\Gamma_i=R_i<I_i=\Lambda_i$, and the identity on the $\varphi,I$ blocks; so $\Theta_n$ is a $1$-chart. The composition of $1$-charts is a $1$-chart (Remark~\ref{bpP:p1:typical}(3)); $\Theta_n(O_0)\subset\mathcal O_n$ (indeed $=\{\Gamma>0\}\cap\mathcal O_n$).
\proofstep{(b) Invariance}\label{stepP:p7:la-inv} $\Theta_n$ is $2\pi$-periodic in $\psi$ and equals $\Theta_n$ in $\varphi$ shifted; precisely $\Theta_n(z+\iota'_k)=\Theta_n(z)+((2\pi k_\varphi,0),0)$ ($k=(k_\varphi,k_\psi)$), and $\Psi^{(n)}$ is invariant under $((2\pi k_\varphi,0),0)$ (Theorem~\ref{bpP:p7:product-chart}(a)). If $\Phi_n(z')=\Phi_n(z)$ then $\Theta_n(z')=\Theta_n(z)+((2\pi k_\varphi,0),0)$ by the same theorem, so $\varphi'=\varphi+2\pi k_\varphi$, $I'=I$, and $\sqrt{2R'}(\sin\psi',\cos\psi')=\sqrt{2R}(\sin\psi,\cos\psi)$ with $R,R'>0$: then $R'=R$ and $\psi'-\psi\in2\pi\Z^n$.
\proofstep{(c) Period one}\label{stepP:p7:la-per} Lemma~\ref{bpP:p1:period-convert}(a): $S_{2\pi}$ is a $2\pi$-chart with $S_{2\pi}(z+\iota_k)=S_{2\pi}(z)+(2\pi k,0)$; compose with (a), (b) (Remark~\ref{bpP:p1:typical}(3): $\Phi_n\circ S_{2\pi}$ is a $2\pi$-chart), and note $S_{2\pi}^{-1}(O_0)+\iota_k=S_{2\pi}^{-1}(O_0+\iota'_k)=S_{2\pi}^{-1}(O_0)$; the injectivity statement (Q2) for $\Phi_n\circ S_{2\pi}$ is that of (b) read through $S_{2\pi}$.
\proofstep{(d)}\label{stepP:p7:la-ham} Theorem~\ref{bpP:p7:product-chart}(b),(c) and $c=1$.
\end{proof}

\subsection{The D25 secular definitions in terms of the chart}\label{bpP:p7:sec-d25}

The frozen D25 definitions (\texttt{EllipsePosition.lean}, \texttt{EllipseMomentum.lean}, \texttt{KeplerMeanLongitude.lean}, \texttt{LongitudeJacobian.lean}, \texttt{SecularEccVector.lean}, \texttt{SecularRoot.lean}, \texttt{SecularAxis.lean}, \texttt{SecularGamma.lean}, \texttt{EllipsePhase.lean}, \texttt{KeplerAction.lean}, \texttt{KeplerAxis.lean}, \texttt{KeplerEnergy.lean}, \texttt{KeplerEccVector.lean}) are stated for $n=2$ and $G=m_0=1$ ($\kappa=1$) in the phase of one chart point $w\in\texttt{Phase}(2+2)$ with
$w_1=(\lambda_1,\lambda_2,\xi_1,\xi_2)$, $w_2=(\Lambda_1,\Lambda_2,\eta_1,\eta_2)$ (so $\texttt{fast}\,w_1=\lambda$, $\texttt{slow}\,w_1=\xi$, $\texttt{fast}\,w_2=\Lambda$, $\texttt{slow}\,w_2=\eta$) and a vector of eccentric longitudes $F\in\R^2$.
(The same statement holds for every $n$ componentwise.)

\begin{proposition}[D25 definitions = the chart in the eccentric longitude]\label{bpP:p7:d25}\label{bpP:p7:eccentric-longitude}
Let $n=2$, $\kappa=1$, $m\in(0,\infty)^2$, $w\in\mathcal O_2$, $i<2$, and write $\Lambda_i=\texttt{fast}\,w_2\,i$, $\Gamma_i=\texttt{secularGamma}\,w\,i$. Then (all formulas for the Lean functions, which are total, are guarded by $\Lambda_i>0$, $\Gamma_i<\Lambda_i$):
\begin{enumerate}[label=(\alph*)]
\item \textup{(vector, root, axis)} $\Gamma_i=\tfrac12(\xi_i^2+\eta_i^2)\ge0$; \Lean{secularRoot}${}=k_i=1-\Gamma_i/\Lambda_i\in(0,1]$; \Lean{secularEccVector}${}=v_i$ of Definition~\ref{bpP:p7:poincare-def} with $\pl{v_i}^2=1-k_i^2$ (eccentricity $e_i=\pl{v_i}$); \Lean{secularAxis}${}=(\Lambda_i/m_i)^2=a_i$.
\item \textup{(Kepler's equation, Jacobian)} $F\mapsto$\Lean{keplerMeanLongitude}$\,w\,i\,F=\Phi(F,v_i)$ is a real-analytic increasing bijection of $\R$ commuting with $F\mapsto F+2\pi$, with derivative \Lean{longitudeJacobian}${}=J(F,v_i)\ge1-e_i>0$. Let $F_i(w)$ be the unique solution of $\Phi(F,v_i)=\lambda_i$ (Proposition~\ref{bpP:p7:regular-kepler}(a); $\lambda_i=\texttt{fast}\,w_1\,i$); $F_i$ is real analytic on $\mathcal O_2$.
\item \textup{(positions and momenta)} For all $F\in\R$ and $i<2$, with $\tilde q_i(F)=$\Lean{ellipsePosition}$\,m\,w\,i\,F$, $\tilde p_i(F)=$\Lean{ellipseMomentum}$\,m\,w\,i\,F$:
$(\tilde q_i(F_i(w)),\tilde p_i(F_i(w)))=\Psi_{m_i,1}(\lambda_i,\xi_i,\Lambda_i,\eta_i)$; i.e.\ \Lean{ellipsePhase}$\,m\,w\,(F_1(w),F_2(w))=\Psi^{(2)}(w)$.
More generally $\tilde q_i(F)$, $\tilde p_i(F)$ for \emph{arbitrary} $F$ are the point of the same osculating Kepler ellipse (same $E$, $\operatorname{cr}$, $v$) at the eccentric longitude $F$, i.e.\ $\Psi(\Phi(F,v_i),\xi_i,\Lambda_i,\eta_i)$.
\item \textup{(Kepler functions of the physical point)} For $z=\Psi^{(2)}(w)$: \Lean{ProgradeElliptic}$\,m\,z$; \Lean{keplerEnergy}$\,m\,z\,i=-m_i^3/(2\Lambda_i^2)$; \Lean{keplerAxis}$\,m\,z\,i=a_i=$ \Lean{secularAxis}$\,m\,w\,i$; \Lean{keplerAction}$\,m\,z\,i=\Lambda_i$; \Lean{planeCross}$(q^{(i)},p^{(i)})=\Lambda_i-\Gamma_i$; \Lean{keplerEccVector}$\,m\,z\,i=v_i=$ \Lean{secularEccVector}$\,w\,i$. Conversely every $z$ with \Lean{ProgradeElliptic}$\,m\,z$ is $\Psi^{(2)}(w)$ for a $w\in\mathcal O_2$, unique modulo $\lambda\mapsto\lambda+2\pi k$.
\item \textup{(elements)} For $\Gamma_i>0$: $(\xi_i,\eta_i)=\rho_ie_{g_i}$, $g_i$ the perihelion argument, $F_i=g_i+u_i$, $u_i$ the eccentric anomaly and $\lambda_i=\ell_i+g_i$ (mean longitude = mean anomaly + perihelion argument).
\end{enumerate}
\end{proposition}
\begin{proof}
\proofstep{(a),(b)}\label{stepP:p7:d25-ab} The Lean definitions are exactly $\Gamma=R(\eta,\xi)$, $k=1-\Gamma/\Lambda$, $v=\sqrt{(1-\Gamma/(2\Lambda))/\Lambda}\,(\xi,\eta)$, $\Phi(F,v)=F-v_0\sin F+v_1\cos F$, $J=1-v\cdot(\cos F,\sin F)$, $a=(\Lambda/m)^2$ (\Lean{planeDot}, \Lean{orbitalUnit}); Lemma~\ref{bpP:p7:o-facts} and Proposition~\ref{bpP:p7:regular-kepler}(a).
\proofstep{(c)}\label{stepP:p7:d25-c} \Lean{ellipsePosition}${}=a(kc-v+\frac{v\cdot c}{1+k}v)$ and \Lean{ellipseMomentum}${}=\frac{m^2}{\Lambda J}(J_0c-\frac{v\cdot c}{1+k}J_0v)$, $c=\texttt{orbitalUnit}\,F$ (note $\texttt{fast}\,w_2\,i\cdot\texttt{longitudeJacobian}=\Lambda J$): these are the displayed $q,p$ of Definition~\ref{bpP:p7:poincare-def} with $\kappa=1$ and $F$ arbitrary. At $F=F_i(w)$ they are $\Psi$. For arbitrary $F$ put $\lambda'=\Phi(F,v_i)$: then $F=F(\lambda',v_i)$ by uniqueness, so the same formulas give $\Psi(\lambda',\xi_i,\Lambda_i,\eta_i)$.
\proofstep{(d)}\label{stepP:p7:d25-d} Theorem~\ref{bpP:p7:poincare}(d): $\Psi(w)\in\mathcal K_{m_i,1}$ blockwise, i.e.\ $m_i>0$, $\pl{q^{(i)}}>0$ (\Lean{planeLength}$>0$), $E<0$ (\Lean{keplerEnergy} is $E_{m_i,1}$), $\operatorname{cr}>0$: this is \Lean{ProgradeElliptic}. The values: $E=-m_i^3/(2\Lambda_i^2)$; \Lean{keplerAxis}${}=-m_i/(2E)=\Lambda_i^2/m_i^2$; \Lean{keplerAction}${}=m_i\sqrt{a}=\Lambda_i$; the LRL identity of Theorem~\ref{bpP:p7:poincare}(d) is \Lean{keplerEccVector}${}=v_i$. Converse and uniqueness: Theorem~\ref{bpP:p7:product-chart}(a).
\proofstep{(e)}\label{stepP:p7:d25-e} Step~\ref{stepP:p7:po-del} and Proposition~\ref{bpP:p7:regular-kepler}(c).
\end{proof}

\begin{corollary}[Angular averages via the eccentric longitude]\label{bpP:p7:average}
Let $w\in\mathcal O_2$ and $\varphi:\R^4\to\R$ be continuous. Then for each fixed $(\xi,\Lambda,\eta)$, as functions of $\lambda_i$ with the other $\lambda_j$ fixed,
\[\frac1{2\pi}\int_0^{2\pi}\varphi\bigl(\Psi^{(2)}(\lambda,\xi,\Lambda,\eta)\bigr)\,d\lambda_i=\frac1{2\pi}\int_0^{2\pi}\varphi\bigl(\texttt{ellipsePhase}\,m\,w\,F\bigr)\Big|_{F_i=s}\,J(s,v_i)\,ds\]
for the unit speed $F_i=s$ and the other $F_j=F_j(w)$ fixed (the integrand is the point of the $i$-th osculating ellipse at eccentric longitude $s$), with $J=$ \Lean{longitudeJacobian}. In particular, with $a_i,e_i$ and (Proposition~\ref{bpP:p7:kepler}(c)):
$\langle1/\pl{q^{(i)}}\rangle_{\lambda_i}=1/a_i$, $\ \langle\pl{p^{(i)}}^2\rangle_{\lambda_i}=m_i^2/a_i$, $\ \langle\pl{q^{(i)}}\rangle_{\lambda_i}=a_i(1+e_i^2/2)$, with $a_i=(\Lambda_i/m_i)^2$ and $e_i^2=\pl{v_i}^2=1-k_i^2$.
\end{corollary}
\begin{proof}
\proofstep{Change of variable}\label{stepP:p7:av-cov} For fixed $v=v_i$, $s\mapsto\Phi(s,v)$ is $C^1$, increasing, with $\Phi(s+2\pi)=\Phi(s)+2\pi$ and $\Phi'=J>0$ (Proposition~\ref{bpP:p7:regular-kepler}(a)); its inverse is $\lambda\mapsto F(\lambda,v)$. Over one period, $\int_{\lambda_*}^{\lambda_*+2\pi}\Xi(F(\lambda,v))\,d\lambda=\int_{F_*}^{F_*+2\pi}\Xi(s)J(s,v)\,ds$ for continuous $2\pi$-periodic $\Xi$ (\ref{bpP:p7:e-cov}), with $\Xi(s)=\varphi(\Psi\text{ at eccentric longitude }s)$ by Proposition~\ref{bpP:p7:d25}(c); periodic integrals do not depend on the starting point. \proofstep{Means}\label{stepP:p7:av-mean} The mean values: the average over $\lambda_i$ over a period equals that over $\ell_i=\lambda_i-g_i$ for $\Gamma_i>0$ (shift by a constant of a periodic function), and Proposition~\ref{bpP:p7:kepler}(c) applies ($\kappa=1$, $m=m_i$, $\Lambda=\Lambda_i$, $C=\Lambda_ik_i$, $a=a_i$); for $\Gamma_i=0$ the circular case of the same item (\Iref{bpI:p7:circular}).
\end{proof}

\subsection{Sources, deviations, errata, forward references}\label{bpP:p7:notes}

\paragraph{Sources read.} Fej04r pp.~46--48 (Hamiltonian (28), the heliocentric coordinates, (31) and Thm.~62, (32)--(33)), Fej07e p.~3 (\texttt{Formule (31)}: factor $\sqrt2$; \texttt{Formules (35--37)}: sign conventions and the sign test). Chi05 \S2 via track~I's Part~7 (not re-read here). The D25 frozen Lean files listed in \S\ref{bpP:p7:sec-d25}.

\paragraph{Totality.} The chart uses $1/\Lambda$, $\sqrt{(1-\Gamma/2\Lambda)/\Lambda}$, $1/(1+k)$, $1/J$: guarded on $\mathcal O$ by $\Lambda>0$, $\Gamma<\Lambda$ (so radicand$>\tfrac1{2\Lambda}>0$, $k\in(0,1]$, $J\ge1-e>0$; Lemma~\ref{bpP:p7:o-facts}). The frozen Lean functions \Lean{secularEccVector}, \Lean{ellipsePosition}, \Lean{ellipseMomentum} use total \Lean{Real.sqrt} and division; on $\mathcal O_2$ nothing is truncated. The inverse of the matrix $k\mathrm{Id}+vv^\top/(1+k)$ is used only in the formula of Theorem~\ref{bpP:p7:poincare}(e), where its eigenvalues are $k,1>0$.

\paragraph{Deviations from the source.} (1) Fej04r (31) and (32) use the reduced mass $\mu_j$ and the corresponding total mass $M_j$ in the Kepler part; the frozen $H_{\rm mov}$ (and Part~6) put $\mu\pl{p_i}^2/(2m_0)$ into the perturbation, so the unperturbed Kepler problem has mass $m_i$ and strength $Gm_0$ and Fej04r's $\Lambda_j$ is replaced by $m_i\sqrt{Gm_0a_i}$ (Definition~\ref{bpP:p7:kepler-data}, Remark~\ref{bpP:p7:same-chart}). (2) The planar chart has no nodal angle: $\theta_j=0$ and the spatial coordinates $(p_j,q_j)$ of (31) are absent. (3) The proof of Thm.~62 is ours (not Poincar\'e/Chenciner).
(4) The D25 secular definitions are stated for $G=m_0=1$; the chart is proved for general $\kappa=Gm_0$ (Theorem~\ref{bpP:p7:poincare}), and Proposition~\ref{bpP:p7:d25} is the case $\kappa=1$.

\paragraph{Suspected errata.} (1) Fej04r (31), p.~47 (original printing as used in Fej07e): ``$r_j=\xi_j+i\eta_j=\sqrt{\Lambda_j}\sqrt{1-\sqrt{1-\varepsilon_j^2}}\ldots$'' without the factor $\sqrt2$ makes the chart only conformally symplectic: already recorded in Fej07e p.~3, reproduced here as Remark~\ref{bpP:p7:normalisation}(1). (2) The orchestrator's brief wrote the regularising change as $(\xi,\eta)=\sqrt{2(\Lambda-G)}(\cos,-\sin)(g)$; with the frozen order $(\lambda,\xi;\Lambda,\eta)$ (position $\xi$) the correct sign is $(\cos g,+\sin g)$ (derived in Step~\ref{stepP:p7:po-symp} and checked numerically in a scratch script); the minus sign belongs to the exchanged-roles convention (Remark~\ref{bpP:p7:normalisation}(2)). Not an error in a pinned source.

\paragraph{Doubts about frozen statements.} None. The frozen D25 definitions agree with the chart (checked above: formulas, and by a numerical scratch check at one sample point that $\Psi^*\Omega=\Omega$, $E=-m^3/2\Lambda^2$, $\operatorname{cr}=\Lambda-\Gamma$, LRL $=v$ at a sample point). \Lean{secularAxis}$=(\Lambda/m)^2$ requires $Gm_0=1$ (D25 normalisation), consistent with Remark~\ref{bpP:p6:units}.

\paragraph{External facts used.} (E1)--(E7) of \S\ref{bpP:p7:conv}, plus \Kref{bp:p1:real-restriction} and the results of Parts~1 and~6 and track~I quoted by label.

\paragraph{Forward references and consumers.} Part~8: Corollary~\ref{bpP:p7:average}, Proposition~\ref{bpP:p7:kepler}(c), Theorem~\ref{bpP:p7:product-chart}(b),(c), Remark~\ref{bpP:p7:same-chart}; Parts~2, 5: Proposition~\ref{bpP:p7:holomorphic}; Parts~12, 13: Lemma~\ref{bpP:p7:lattice}, Lemma~\ref{bpP:p7:tube}, Theorem~\ref{bpP:p7:product-chart}; Part~14: Proposition~\ref{bpP:p7:d25}, Corollary~\ref{bpP:p7:average}.
Needed from other parts and not yet written (forward): none beyond the labels of Parts~1, 6 and of track~I Part~7 already in the repository.

\paragraph{Label index.} \texttt{bpP:p7:kepler-data}, \texttt{range}, \texttt{kepler}, \texttt{regular-kepler}, \texttt{poincare-def}, \texttt{o-facts}, \texttt{poincare}, \texttt{holomorphic}, \texttt{normalisation}, \texttt{product-def}, \texttt{product-chart}, \texttt{same-chart}, \texttt{tube}, \texttt{lattice}, \texttt{d25} (= \texttt{eccentric-longitude}), \texttt{average}.
\par\endgroup
